\documentclass[10pt,a4paper]{amsart}
\usepackage[margin=19mm]{geometry}
\usepackage{amsmath,amssymb,mathtools}
\usepackage[scr=rsfs,cal=euler]{mathalfa}
\usepackage{xfrac}
\usepackage[unicode,hidelinks,bookmarks=false]{hyperref}
\newcommand{\ie}{i.e.\ }
\newcommand{\cf}{cf.\ }
\newcommand{\resp}{respectively\ }
\newcommand{\Subsection}{\subsection}
\providecommand{\nobreakdash}{\nobreak\hbox{-}}
\setlength{\emergencystretch}{2em}
\multlinegap0pt
\usepackage{bm}
\usepackage{bbm}
\usepackage{dsfont}
\usepackage{xspace}
\usepackage{cancel}
\usepackage{mathtools}
\usepackage{amsthm}
\makeatletter
\@ifundefined{theorem}{\newtheorem{theorem}{Theorem}}{}
\@ifundefined{corollary}{\newtheorem{corollary}[theorem]{Corollary}}{}
\@ifundefined{proposition}{\newtheorem{proposition}[theorem]{Proposition}}{}
\makeatother
\def\Rn{\mathbb R}
\def\x{\bar {x}}\def\y{\bar {y}}
\title[Affine logic with the integration operator]{Affine logic with the integration operator}
\author{Seyed-Mohammad Bagheri}
\date{Source: arXiv:2602.19056v1 (CC BY 4.0)}
\begin{document}
\maketitle

\noindent\emph{Source and licence.} Affine logic with the integration operator. Version arXiv:2602.19056v1 (CC BY 4.0), distributed under the Creative Commons Attribution 4.0 International licence, as recorded in the arXiv licence metadata for this exact version and shown on the abstract page https://arxiv.org/abs/2602.19056v1. Full rights notice: licenses/arxiv-2602.19056-CC-BY-4.0.txt.\par\smallskip

\noindent\textit{Editorial scope.} Counted unit: the Proposition whose source label is \texttt{model construction} in the article (source0.tex lines 634--636, source label \texttt{model construction}), delivered together with the article's own complete original proof of it (source0.tex lines 637--668, one \texttt{proof} environment of 32 lines). No mathematics is written, repaired or re-derived: statement and proof are the article's own text line for line. Required context retained uncounted: KR (lines 250--254); exist1 (lines 278--282). Every definition, display and label of those lines is retained unchanged. Omitted from this package and not redistributed: 1--249, 255--277, 283--633, 669--1277. Nothing inside a retained excerpt is omitted or altered: every retained body is delivered exactly as the article's lines are written, so no omission receipt of any kind accompanies this package. Citations: the retained excerpts carry no citation command at all, so no bibliography entry of the article is transcribed and no reference is redistributed. Reference adaptation: none was needed and none was made; every reference inside the delivered unit resolves inside it, so no unresolved reference or citation is produced. Printed slips kept literal: no printed word, symbol, hypothesis or number of the counted statement or of its proof was altered. Layout and class: the article's own class is \texttt{article} with packages including amsmath, amsthm, amssymb, amsfonts, mathrsfs, hyperref; the wrapper is the lane's pinned \texttt{amsart} editorial wrapper at 10pt on A4, which restates the theorem-like environments the retained text needs and adds the article's own macro definitions that the retained text needs (\texttt{\textbackslash Rn}, \texttt{\textbackslash x}), with extra wrapper packages (bm, bbm, dsfont, xspace, cancel, mathtools). The delivered numbering is the unit's own. Assets: no retained span contains a figure environment, a graphics-inclusion command, a PDF-page inclusion, a table or a TikZ picture; no image file is copied into this package, the package declares no asset dependency and no image file is redistributed with it. Licence: version arXiv:2602.19056v1 of this article is distributed under the Creative Commons Attribution 4.0 International licence, as recorded in the arXiv licence metadata for this exact version and shown on the abstract page https://arxiv.org/abs/2602.19056v1. A full rights notice is delivered at licenses/arxiv-2602.19056-CC-BY-4.0.txt. Characters: every retained excerpt is pure ASCII, so the delivered engine, XeLaTeX with its default Latin Modern text font, reports no missing character.
\begin{corollary}\label{KR}
Let $X$ be a metric space and $G\subseteq{\mathbf C}_b(X)$ be a majorizing vector subspace.
Then, for every positive linear $\Lambda:G\rightarrow\Rn$, there is a regular Borel charge $\mu$ on $X$
such that $\mu(X)=\Lambda(1)$ and $$\Lambda(f)=\int fd\mu\hspace{18mm} \forall f\in G.$$
\end{corollary}

\begin{proposition}\label{exist1}
For each $L$-prestructure $M$ there is a $L$-structure $K$ and a dense map
$\pi:M\rightarrow K$ such that for every formula $\phi(\x)$ and $a_1,...,a_k\in M$
$$\phi^{M}(a_1,...,a_k)=\phi^{K}(\pi a_1,...,\pi a_k).$$
\end{proposition}

\begin{proposition}\label{model construction}
Let $T$ be a maximal affinely satisfiable theory in $L$ with the witness and mean value properties. Then $T$ has a model.
\end{proposition}

\begin{proof}
By maximality, for every sentence $\phi$ there is a unique $r$ such that $\phi=r\in T$.
We denote this $r$ by $\phi^T$.
Let $C$ be the set of constant symbols of $L$ and define a pseudometric on $M$
(which we denote it by $\rho$) by setting $\rho(c,e)=(\rho(c,e))^T$.
Let $M$ be the resulting quotient metric space where its metric is denoted by $\rho$ again.
The equivalence class of $c$ is denoted by $\hat c$. Define a metric structure on $M$ by setting:

- $c^M=\hat c$

- $R^M(\hat a_1,...,\hat a_n)=(R(a_1,...,a_n))^T$

- $F^M(\hat a_1,...,\hat a_n)=\hat b$ \ \ if \ \ $\rho(F(a_1,...,a_n),b)=0\in T$.
\vspace{1mm}

Note that $R^M$ and $F^M$ are well-defined. For every $L$-formula $\phi(x)$,
consider the function $\phi^T:M\rightarrow\Rn$ defined by $\hat c\mapsto(\phi(c))^T$.
The set of such functions is a majorizing subset of $\mathbf{C}_b(M)$.
It is clear that the functional $$\Lambda:\phi^T(x)\mapsto\Big(\int\phi(x)dx\Big)^T$$ is linear.
On the other hand, if $0\leqslant\phi^T(x)$ as a function, then $0\leqslant\phi(c)\in T$ for every $c$.
Hence, by the mean value property, for some $c$ the conditions $$0\leqslant\phi(c)=\int\phi(x)dx$$ belong to $T$.
This shows that $\Lambda$ is positive.
Now, using Corollary \ref{KR}, we obtain a charged $L$-structure $(M,\mu)$ such that for every $\phi(x)$
$$\int\phi^T(x)d\mu=\Big(\int\phi(x)dx\Big)^T.$$
Finally, we must prove that for every $L$-formula $\phi(\x)$ and $a_1,...,a_n\in C$,
$$\phi^M(\hat a_1,...,\hat a_n)=(\phi(a_1,...,a_n))^T.$$
We consider the integral case. Assume the claim is proved for $\phi(x)$ so that
for every $a\in C$, $\phi^M(\hat a)=(\phi(a))^T$.
Then, by the construction of $\mu$ and the induction hypothesis,
$$\Big(\int\phi(x)dx\Big)^T=\int\phi^T(x)d\mu=\int\phi^M(x)d\mu.$$
To complete the proof, use Proposition \ref{exist1} to find a complete model of $T$. %The second part is proved similarly.
\end{proof}

\end{document}
